\section{The intrinsic filtration and its graded map}\label{sec:filtration}

Let \(k\) have characteristic zero, let \(L,M\) be finite-dimensional,
and suppose that \(L\) is solvable and that
\(\Phi_{\mathrm{raw}}:U(L)\xrightarrow{\sim}U(M)\) is a \(k\)-algebra isomorphism.
The preceding results imply that \(M\) is solvable. Write
\[
 N=N(L),\quad N'=N(M),\quad
 Q=L/N,\quad Q'=M/N',
 \qquad I=U(L)N,\quad I'=U(M)N'.
\]
All ideals generated in an enveloping algebra are two-sided. We identify
Lie elements with their canonical images by PBW.

We first fix the normalization from \cref{lem:augmentation}. If
\(\chi(x)=\varepsilon_M(\Phi_{\mathrm{raw}}(x))\), then
\[
 \Phi=\Phi_{\mathrm{raw}}\circ\tau_{-\chi},\qquad
 \tau_{-\chi}(x)=x-\chi(x)1
\]
preserves augmentation. Both \(\Phi\) and \(\Phi^{-1}\) will always refer
to these specified maps. This is essential: the automorphism
\(x\mapsto x+1\) of \(k[x]\) does not preserve \((x)\).
There is no claim that the unnormalized \(\Phi_{\mathrm{raw}}\) preserves \(I\).

\subsection{The metabelian map and its linear terms}

Put \(E=U(L)[N,N]\), and define \(E'\) similarly.
By \cref{prop:intrinsic-ideal}, these ideals are intrinsically generated
by commutators of locally nilpotent inner elements. Thus even
\(\Phi_{\mathrm{raw}}\) carries \(E\) onto \(E'\). In particular, \(\Phi\) induces
an augmentation-preserving isomorphism
\[
 \overline\Phi:U(g)\xrightarrow{\sim}U(g'),
 \qquad g=L/[N,N],\quad
 g'=M/[N',N'].
\]
The abelian ideals \(V=N/[N,N]\) and \(V'=N'/[N',N']\)
embed as polynomial subalgebras \(S(V)\) and \(S(V')\).
Choose linear sections \(s:Q\to g\) and
\(s':Q'\to g'\). They are not assumed to be Lie maps.
The intrinsic descriptions in \cref{prop:LF-LN} give
\[
 \LN(U(g))=S(V),\qquad
 \LF(U(g))=s(Q)+S(V).
\]
The analogous equalities hold for \(g'\).
Consequently the actual map has the form
\begin{equation}\label{eq:metabelian-map}
 \overline\Phi(v)=F(v),\qquad
 \overline\Phi(s(q))=s'(Tq)+b_q,
\end{equation}
where \(F:S(V)\xrightarrow{\sim}S(V')\) is an algebra isomorphism,
\(T:Q\xrightarrow{\sim}Q'\) is linear, and \(q\mapsto b_q\in S(V')\)
is linear. Indeed, \(T\) is the induced map on
\(\LF/\LN\), and PBW identifies this vector-space quotient with \(Q\).
The decomposition is unique since \(s'(Q')\cap S(V')=0\).
Augmentation preservation says that \(F(v)\) and \(b_q\) have zero
constant term.

Let \(P:V\to V'\) be the linear term of \(F\) at the origin, and let
\(B(q)\in V'\) be the linear term of \(b_q\).
The inverse of \(F\) also fixes the origin, so comparison of linear
terms in both compositions proves that \(P\) is invertible.
To keep track of the extension cocycle, write
\[
 [s(q),v]=A_qv,\qquad [s(q),s(r)]=c(q,r)\in V,
\]
and use primes for the corresponding target data. Applying
\(\overline\Phi\) to these commutators and taking linear terms gives
\begin{align}
 PA_q&=A'_{Tq}P,\label{eq:linear-action}\\
 Pc(q,r)&=c'(Tq,Tr)+A'_{Tq}B(r)-A'_{Tr}B(q).
 \label{eq:linear-cocycle}
\end{align}
Here each \(A'_{Tq}\) extends as a derivation preserving ordinary
polynomial degree, so higher terms of \(F\) or \(b_q\) cannot contribute.
Thus the linear map
\(v+s(q)\mapsto Pv+s'(Tq)+B(q)\) is a Lie isomorphism of the
metabelian quotients. In particular, the argument retains the
possibly nonzero cocycle rather than assuming a split extension.

\begin{proposition}\label{prop:filtered-linearization}
The normalized map satisfies \(\Phi(I^d)=(I')^d\) for every
\(d\geq0\). On \(U(L)/I=S(Q)\), its induced map is exactly \(S(T)\).
Under the natural right-module identifications
\[
 I/I^2\simeq V\otimes_k S(Q),\qquad
 I'/(I')^2\simeq V'\otimes_k S(Q'),
\]
its first graded map is \(P\otimes S(T)\).
\end{proposition}

\begin{proof}
The ideal \(I/E\) in \(U(g)\) is generated by the augmentation
ideal \((V)\subset S(V)\). Since \(F\) and \(F^{-1}\) fix the origin,
they carry \((V)\) and \((V')\) onto one another.
Lifting through \(E,E'\) yields \(\Phi(I)=I'\), hence equality of all
their powers. In \eqref{eq:metabelian-map}, the terms \(b_q\)
vanish modulo \(I'\); the quotient map is therefore \(S(T)\).

The PBW description of the first quotients follows from
\cref{lem:ideal-powers} below: order the nilradical variables before
the section variables. The monomials of weight one are a single
representative of \(V\), followed by an arbitrary monomial in \(Q\).
Their classes give the displayed right-module identifications.
The relation \(qv-vq=A_qv\) describes the left action.

Since \([N',N']\subset(I')^2\), we have \(E'\subset(I')^2\).
For \(n\in N\), its image on the first layer may therefore be computed
from \(F(\bar n)\) in the metabelian quotient. This polynomial contains
no \(Q'\)-variables. Its terms of degree at least two in \(V'\) vanish
modulo \((I')^2\), leaving precisely \(P(\bar n)\).
Right multiplication and the already identified quotient map prove
the asserted formula on the whole module. In particular the actual
constant subspace \(V\otimes1\) maps onto \(V'\otimes1\);
an equality of module ranks would not give this conclusion.

Explicitly, on this first quotient the left action is
\[
 q\cdot(v\otimes h)=A_qv\otimes h+v\otimes qh.
\]
Its transport is consistent with \eqref{eq:linear-action}.
Although \(\Phi(s(q))-s'(Tq)\) may contain terms from \(I'\),
commuting any such term with an element of \(I'\) gives an element
of \((I')^2\). Thus these terms introduce no correction to the
first-layer action. This also explains why the term \(B(q)\),
which is required to retain the metabelian extension cocycle,
does not enter the degree-zero map of the descending associated
graded algebra.
\end{proof}

\subsection{Ideal powers and weighted PBW}

Define a descending Lie filtration by
\[
 F^0L=L,\qquad F^dL=\gamma_dN\quad(d\geq1),\qquad
 \gamma_1N=N,\quad \gamma_{d+1}N=[N,\gamma_dN].
\]
Choose a basis \(e_i\) adapted to this finite flag, with weight \(w_i=d\)
on a representative of \(F^dL/F^{d+1}L\). Thus the complement of \(N\)
has weight zero. Let \(W_d\) be the span of ordered PBW monomials
of total weight at least \(d\).

\begin{lemma}\label{lem:ideal-powers}
For every \(d\geq0\), \(I^d=W_d\). For \(d\geq1\),
\[
 L\cap I^d=\gamma_dN.
\]
\end{lemma}

\begin{proof}
The lower central series satisfies
\[
 [\gamma_iN,\gamma_jN]\subset\gamma_{i+j}N,\qquad
 [L,\gamma_jN]\subset\gamma_jN.
\]
Accordingly, every basis vector occurring in \([e_i,e_j]\)
has weight at least \(w_i+w_j\).
In a PBW reordering step \(e_je_i=e_ie_j+[e_j,e_i]\),
no resulting term has smaller weight. Repeated reordering shows
that \(W_dW_e\subset W_{d+e}\). Since \(W_0=U(L)\), \(W_1\)
is a two-sided ideal containing \(N\). Hence \(I^d\subset W_d\).

Conversely, induction using
\([I^r,I^s]\subset I^{r+s}\) gives \(\gamma_jN\subset I^j\).
Every basis vector \(e_i\) belongs to \(I^{w_i}\), including
the weight-zero case. A PBW monomial of total weight \(w\)
therefore lies in \(I^w\), and in \(I^d\) when \(w\geq d\).
This proves the reverse inclusion.
Finally, a linear Lie element has only the degree-one PBW coordinates
of its basis expansion. Its belonging to \(W_d\) is exactly its
belonging to the span of the \(e_i\) with \(w_i\geq d\),
which is \(\gamma_dN\).
\end{proof}

The use of zero weights is deliberate. Weight pieces need not be
finite-dimensional, but every PBW expansion has finite support.
In particular \(\bigcap_d I^d=0\), and no finite-dimensionality
of the individual graded pieces has been assumed.
Indeed, for a nonzero element choose any monomial with nonzero
PBW coefficient. Membership in every \(W_d\) would require that
one fixed monomial to have arbitrarily large weight.
The basis and ordering only describe these subspaces: their equality
with the intrinsically specified powers \(I^d\) makes the resulting
filtration independent of both choices.

\subsection{The prescribed graded Lie identification}

The quotient \(Q\) acts canonically on \(\gr_\gamma N\).
Changing a lift of \(q\) by \(n\in N\) changes its action on
\(\gamma_dN\) only in \(\gamma_{d+1}N\).
The commutator of two lifted actions is inner by an element of \(N\)
and hence acts trivially on each graded piece. Thus
\[
 \gr_F L=Q\ltimes\gr_\gamma N,
\]
with \(Q\) in degree zero. Sending a Lie class to its ideal-adic
symbol defines a graded algebra isomorphism
\begin{equation}\label{eq:graded-pbw}
 U(\gr_F L)\xrightarrow{\sim}\gr_I U(L).
\end{equation}
Indeed, the filtration bracket identities make this a homomorphism,
and \cref{lem:ideal-powers} identifies the ordered PBW monomials
of each weight as bases on both sides.

\begin{proposition}\label{prop:graded-mark}
Under \eqref{eq:graded-pbw} and its target analogue, the actual map
\(\gr_I\Phi\) is \(U(\varphi_0)\) for a graded Lie isomorphism
\[
 \varphi_0:Q\ltimes\gr_\gamma N\xrightarrow{\sim}
       Q'\ltimes\gr_\gamma N'.
\]
There is a filtered linear isomorphism \(f:L\to M\) inducing
\(\varphi_0\), such that, for \(x\in F^dL\),
\[
 \Phi(x)-f(x)\in(I')^{d+1}.
\]
Its inverse satisfies the reciprocal congruences.
\end{proposition}

\begin{proof}
The Lie algebra \(\gr_\gamma N\) is generated by its degree-one
subspace \(V\). This follows inductively from
\(\gamma_{d+1}N=[N,\gamma_dN]\): replacing an input from \(N\)
by an element of \(\gamma_2N\) changes the bracket only in the next
layer. By \cref{prop:filtered-linearization}, \(\gr_I\Phi\)
carries \(V\) onto \(V'\). Preservation of associative commutators
therefore carries their generated Lie subalgebras onto one another.
On the degree-zero linear subspace its restriction is \(T\).
Together these restrictions give the stated graded Lie isomorphism.
The enveloping algebras are generated by these Lie subspaces, so
equality on them proves \(\gr_I\Phi=U(\varphi_0)\) on the entire
graded algebra.
In particular, the restriction to \(Q\) is the map \(T\) extracted
from the given metabelian quotient map, and the restriction to \(V\)
is its actual derivative \(P\). No preservation of a coproduct is
part of the hypotheses. What singles out the relevant Lie subspace
inside the graded associative algebra is the constant first layer,
together with generation by commutators, rather than any additional
Hopf-algebra structure imposed on the original isomorphism.

For completeness, the linear lift can be chosen while retaining
this exact map. Choose an adapted basis of \(L\), and choose in \(M\)
a representative of \(\varphi_0([e_i])\) in degree \(w_i\).
Extend these choices linearly to \(f\). The equality of symbols gives
\[
 f(F^dL)\subset F^dM,\qquad
 \Phi(x)-f(x)\in(I')^{d+1}\quad(x\in F^dL).
\]
Apply the same construction to \(\Phi^{-1}\), obtaining \(g\).
Ideal-power preservation and \cref{lem:ideal-powers} imply
\[
 (gf-\Id)(F^dL)\subset F^{d+1}L,\qquad
 (fg-\Id)(F^dM)\subset F^{d+1}M.
\]
Both error operators are nilpotent because the Lie filtrations
terminate. Hence both compositions are invertible, so \(f\) is an
isomorphism. Its inverse is filtered: for a nonzero vector of highest
filtration degree, invertibility of the associated graded map
prevents its image from entering the next layer.
Applying \(\Phi^{-1}\) to the forward congruence now proves the
reciprocal one for \(f^{-1}\).
\end{proof}

In particular, with \(e_i'=f(e_i)\), the two adapted bases have
the same weights and satisfy
\begin{equation}\label{eq:matched-leading-bases}
 \Phi(e_i)-e_i'\in(I')^{w_i+1},\qquad
 \Phi^{-1}(e_i')-e_i\in I^{w_i+1}.
\end{equation}
These are the fixed leading data for the Rees construction in the
next section. Its zero-fiber map is consequently the prescribed
\(\varphi_0\), and its entire enveloping map is \(U(\varphi_0)\).
One can also check the graded Lie assertion directly from the leading
congruences. For \(x\in F^dL\) and \(z\in F^eL\), write
\(u=\Phi(x)-f(x)\) and \(v=\Phi(z)-f(z)\). Then
\[
 \Phi([x,z])-[f(x),f(z)]
   =[u,\Phi(z)]+[f(x),v]\in(I')^{d+e+1}.
\]
The leading congruence for \([x,z]\in F^{d+e}L\), followed by
the intersection statement of \cref{lem:ideal-powers}, therefore gives
\[
 f([x,z])-[f(x),f(z)]\in F^{d+e+1}M.
\]
This is bracket preservation on the associated graded algebra,
with precisely the symbols supplied by \(\Phi\).
If another choice of representatives produces a lift \(\widetilde f\),
the same intersection statement gives
\((f-\widetilde f)(F^dL)\subset F^{d+1}M\).
Hence the graded identification is independent of the chosen lift.
The individual lifts serve only to choose compatible finite bases
for later constructions; they do not replace the marked graded map
by a different isomorphism.

Neither the original brackets nor the embedded subalgebra \(U(N)\)
have been asserted to be preserved by the chosen lift \(f\) or by
\(\Phi\), respectively.
